% !TEX root = ../main.tex

\section{Hardvérová časť}

Hardvérová časť predstavuje realizáciu technických požiadaviek uvedených v kapitole 2.2. 
V tejto kapitole sú opísané použité komponenty, hydraulický okruh, 
Peltierova chladiaca zostava, výkonové zapojenie, napájanie a spôsob merania teploty.

Pri návrhu bolo potrebné zohľadniť najmä prúdové zaťaženie výkonových prvkov, 
stabilné napájanie, účinné chladenie teplej strany Peltierovho článku a 
bezpečné oddelenie elektroniky od vodného okruhu.

\subsection{Použité komponenty}

V projekte boli použité komponenty potrebné na vytvorenie hydraulického okruhu, 
merania teploty, výkonového riadenia, napájania a lokálneho zobrazovania údajov. 
Prehľad hlavných použitých komponentov je uvedený v tabuľke \ref{tab:pouzite_komponenty}.

{\small
\begin{longtable}{|>{\RaggedRight\arraybackslash}p{0.30\textwidth}|>{\centering\arraybackslash}p{0.12\textwidth}|>{\RaggedRight\arraybackslash}p{0.48\textwidth}|}
\hline
\textbf{Komponent} & \textbf{Počet} & \textbf{Úloha v systéme} \\
\hline
ESP32 & 1 ks & Riadiaca jednotka systému, čítanie údajov zo snímača, generovanie PWM signálov a komunikácia so serverom. \\
\hline
Peltierova zostava & 1 ks & Aktívne chladenie kvapaliny pomocou Peltierovho článku. \\
\hline
Peltierov článok & 1 ks & Termoelektrický prvok vytvárajúci rozdiel teplôt medzi studenou a teplou stranou. \\
\hline
Chladič s ventilátorom & 1 ks & Odvod tepla z teplej strany Peltierovho článku. \\
\hline
Vodná pumpa 12~V & 1 ks & Cirkulácia kvapaliny v uzavretom hydraulickom okruhu. \\
\hline
Hadičky pre vodné pumpy & podľa potreby & Prepojenie nádržky, pumpy a výmenníka tepla. \\
\hline
Výmenník tepla & 1 ks & Prenos tepla medzi kvapalinou a studenou stranou Peltierovho článku. \\
\hline
Vyhrievacie teleso 40~W & 1 ks & Simulácia tepelnej záťaže alebo poruchového stavu. \\
\hline
Výkonový rezistor 330~$\Omega$ / 10~W & 1 ks & Pomocný výkonový prvok alebo záťaž podľa aktuálneho zapojenia. \\
\hline
H-mostík L298N & 1 ks & Výkonové riadenie Peltierovho článku alebo motorických záťaží podľa zapojenia. \\
\hline
PWM MOSFET modul & 1 ks & Spínanie alebo PWM riadenie výkonovej záťaže. \\
\hline
Step-down menič Mini-360 MP2307 & 1 ks & Zníženie napätia pre vybrané časti zapojenia. \\
\hline
Step-down menič s USB výstupom 5~V / 5~A & 1 ks & Vytvorenie stabilizovanej 5~V vetvy pre riadiacu elektroniku alebo periférie. \\
\hline
LCD displej 16x2 s I2C modulom & 1 ks & Lokálne zobrazovanie základných údajov systému. \\
\hline
Digitálny snímač teploty DS18B20 & 1 ks & Meranie teploty kvapaliny. \\
\hline
Rezistor 4,7~k$\Omega$ & 1 ks & Pull-up rezistor pre dátovú linku snímača DS18B20. \\
\hline
Svorkovnice & 6 ks & Pripojenie napájania, výkonových záťaží a externých vodičov. \\
\hline
Kondenzátor 1000~$\mu$F & 1 ks & Stabilizácia napájania pri spínaní výkonových prvkov. \\
\hline
Kondenzátor 470~$\mu$F & 1 ks & Dodatočné vyhladenie napájacej vetvy. \\
\hline
DC jack / napájací konektor & 1 ks & Pripojenie externého napájacieho zdroja. \\
\hline
Prepojovacie vodiče a konektory & podľa potreby & Elektrické prepojenie jednotlivých častí systému. \\
\hline
\caption{Zoznam hlavných použitých komponentov}
\label{tab:pouzite_komponenty}
\end{longtable}
}

Použité komponenty boli zvolené tak, aby bolo možné vytvoriť funkčný laboratórny model s 
možnosťou merania, regulácie a testovania. Pri výkonových prvkoch bolo potrebné zohľadniť 
ich prúdové zaťaženie, zahrievanie a potrebu spoľahlivého chladenia.

\subsection{Mechanická konštrukcia a hydraulický okruh}

Mechanická časť systému je založená na uzavretom hydraulickom okruhu. Okruh pozostáva z nádržky, 
vodnej pumpy, hadičiek a výmenníka tepla, ktorý je tepelne spojený so studenou stranou 
Peltierovho článku. Úlohou okruhu je zabezpečiť cirkuláciu kvapaliny cez výmenník tak, 
aby bolo možné meniť jej teplotu a sledovať odozvu systému.

Základný smer prúdenia kvapaliny je:

\begin{center}
Nádržka $\rightarrow$ vodná pumpa $\rightarrow$ výmenník tepla pri Peltierovom článku $\rightarrow$ návrat do nádržky
\end{center}

Kvapalina je z nádržky nasávaná pumpou a cez hadičky vedená do výmenníka tepla. Pri zapnutí 
Peltierovho článku sa jeho studená strana ochladzuje a cez výmenník odoberá teplo kvapaline. 
Ochladená kvapalina sa následne vracia späť do nádržky.

Pri mechanickej realizácii bolo potrebné dbať na tesnosť spojov, správne vedenie hadičiek a 
zabránenie úniku kvapaliny v blízkosti elektroniky. Dôležitý je aj dobrý tepelný kontakt medzi 
Peltierovým článkom, výmenníkom a chladičom.

\subsection{Elektronické zapojenie a výkonové riadenie}

Elektronická časť systému prepája mikrokontrolér ESP32, snímač teploty, LCD displej, 
výkonové moduly a napájacie vetvy. ESP32 pracuje ako riadiaca jednotka, ktorá číta teplotu, 
prijíma riadiace hodnoty zo servera a generuje PWM signály pre výkonové členy.

Výkonové prvky, ako Peltierov článok, vodná pumpa, ventilátor a vyhrievacie teleso, 
nie je možné napájať priamo z pinov mikrokontroléra. Preto sú riadené cez výkonové moduly, 
napríklad H-mostík L298N alebo MOSFET modul. Tieto moduly umožňujú spínať alebo regulovať 
záťaže s vyšším prúdovým odberom pomocou nízkoúrovňového riadiaceho signálu z ESP32.

H-mostík umožňuje riadenie výkonovej záťaže a v prípade potreby aj zmenu polarity. Pri 
Peltierovom článku môže zmena polarity zmeniť smer tepelného toku. MOSFET modul je vhodný 
najmä na spínanie alebo PWM riadenie pumpy a vyhrievacieho telesa.

Pri návrhu zapojenia bolo potrebné zabezpečiť spoločnú zem výkonovej a riadiacej časti, 
aby mali PWM signály správnu referenčnú úroveň. Zároveň je vhodné viesť výkonové vodiče 
oddelene od signálových vodičov, aby sa znížil vplyv rušenia na snímač teploty a komunikáciu s perifériami.

\begin{figure}[H]
    \centering
    \safeimage{obrazky/Schema_zapojenia1.png}{0.95\textwidth}{Detailná schéma elektrického zapojenia laboratórneho modelu}
    \label{fig:schema_zapojenia}
\end{figure}

\subsection{Napájanie, návrh dosky a stabilizácia zapojenia}

Napájanie systému je rozdelené na výkonovú a riadiacu časť. Výkonová časť pracuje najmä s 
napätím 12~V, ktoré slúži na napájanie Peltierovho článku, vodnej pumpy, ventilátora a 
vyhrievacieho telesa. Riadiaca časť potrebuje stabilizované nižšie napätie pre ESP32, 
LCD displej a ďalšie logické obvody.

Na vytvorenie nižších napäťových vetiev sú použité step-down meniče. Mini-360 MP2307 slúži 
na zníženie vstupného napätia pre vybrané časti obvodu. Step-down menič s USB výstupom 
5~V / 5~A zabezpečuje stabilizovanú 5~V vetvu pre riadiacu elektroniku alebo periférie.

Pri návrhu dosky bolo potrebné oddeliť výkonové vetvy od signálových vodičov. Výkonové vodiče 
pre Peltierov článok, pumpu a vyhrievacie teleso musia byť dimenzované na vyšší prúd. 
Signálové vodiče pre snímač teploty, I2C displej a PWM riadenie môžu byť vedené menšími 
prierezmi, ale mali by byť vedené tak, aby boli čo najmenej ovplyvnené rušením z výkonovej časti.

V zapojení sú použité kondenzátory 1000~$\mu$F a 470~$\mu$F. Ich úlohou je vyhladzovať 
napájacie napätie a kompenzovať krátkodobé poklesy alebo špičky, ktoré môžu vzniknúť 
pri spínaní pumpy, Peltierovho článku alebo vyhrievacieho telesa. Použitie svorkovníc 
zjednodušuje pripájanie výkonových záťaží a zvyšuje prehľadnosť zapojenia.

\begin{figure}[H]
    \centering
    \safeimage{obrazky/plosak_spodok.png}{0.80\textwidth}{Návrh rozloženia výkonovej časti zapojenia}
    \label{fig:plosak_v1}
\end{figure}

\subsection{Meranie teploty a lokálne zobrazenie}

Teplota kvapaliny je meraná digitálnym snímačom DS18B20. Snímač je vhodný na meranie teploty 
kvapaliny, pretože je dostupný aj vo vodotesnom puzdre a komunikuje digitálne. Pri zapojení 
sa používa dátová linka, napájanie, zem a pull-up rezistor 4,7~k$\Omega$.

Snímač je umiestnený v nádržke alebo v časti okruhu, ktorá reprezentuje aktuálnu teplotu 
kvapaliny. Nameraná hodnota slúži ako hlavná spätnoväzbová veličina systému a používa sa 
na výpočet regulačnej odchýlky.

Lokálne zobrazovanie údajov je riešené pomocou LCD displeja 16x2 s I2C modulom. Displej 
umožňuje zobraziť základné informácie, napríklad aktuálnu teplotu, požadovanú teplotu, 
stav systému, regulačný mód alebo hodnoty PWM. Vďaka I2C rozhraniu stačia na komunikáciu 
iba dve dátové vodiče.

\subsection{Oživenie a vývoj zapojenia počas projektu}

Oživenie hardvéru prebiehalo postupne. Najskôr bolo potrebné overiť napájanie riadiacej časti, 
funkciu ESP32, snímača teploty a LCD displeja. Až následne bolo vhodné pripájať 
výkonové vetvy a testovať pumpu, Peltierov článok, ventilátor a vyhrievacie teleso.

Pri hydraulickej časti bolo potrebné skontrolovať naplnenie okruhu, správnu cirkuláciu 
kvapaliny a tesnosť spojov. Pred zapnutím Peltierovho článku bolo nutné overiť funkčnosť 
chladiča a ventilátora na jeho teplej strane. Počas prvých testov bolo vhodné sledovať 
aj teplotu výkonových modulov, vodičov a svorkovníc.

Hardvérové zapojenie sa počas projektu postupne upravovalo. Prvé verzie slúžili najmä na 
overenie základného princípu zapojenia komponentov, napájacích vetiev a výkonových modulov. 
Neskôr sa návrh upravoval tak, aby bolo zapojenie prehľadnejšie, bezpečnejšie a vhodnejšie na 
praktické testovanie.

Do dokumentácie sú zaradené jednotlivé verzie schém a návrhov zapojenia. Porovnaním starších a 
novších verzií je možné vidieť postupné sprehľadnenie zapojenia a úpravy vyplývajúce z 
praktického testovania.

\begin{figure}[H]
    \centering
    \safeimage{obrazky/fyzicka_realizacia.png}{0.85\textwidth}{Fyzická realizácia laboratórneho modelu}
    \label{fig:fyzicka_realizacia}
\end{figure}
